% Partie : intuitionistic-logic
% Chapitre : introduction
% Section : constructive-reasoning

\documentclass[../../../include/open-logic-section]{subfiles}

\begin{document}

\olfileid[fr]{int}{int}{cr}

\olsection{Raisonnement constructif}

À la différence des extensions de la logique classique par des opérateurs
modaux ou des quantificateurs du second ordre, la logique intuitionniste
est « non classique » en ce qu'elle restreint la logique classique.
Cette dernière est \emph{non constructive} à plusieurs égards. La
logique intuitionniste vise à rendre compte d'un raisonnement plus
« constructif », caractéristique de certaines formes de mathématiques
constructives. Les exemples suivants en éclairent les motivations.

Supposons que quelqu'un affirme avoir déterminé un entier naturel~$n$
tel que, si $n$ est pair, l'hypothèse de Riemann est vraie et, si $n$
est impair, elle est fausse. Excellente nouvelle !{} La vérité de
l'hypothèse de Riemann est l'un des grands problèmes ouverts des
mathématiques. Il semble qu'on ait ramené ce problème à un calcul :
déterminer si un nombre précis est pair ou impair.

Quelle est cette valeur miraculeuse de~$n$ ? Voici comment on la décrit :
$n$ est l'entier naturel égal à $2$ si l'hypothèse de Riemann est vraie,
et à $3$ sinon.

Furieux, vous réclamez le remboursement. Du point de vue classique,
cette description détermine effectivement une unique valeur de $n$ ;
mais ce que vous souhaitez, c'est une valeur de $n$ donnée
\emph{explicitement}.

Prenons un autre exemple, peut-être moins artificiel. Nous savons
qu'élever un nombre irrationnel à une puissance rationnelle peut donner
un résultat rationnel : par exemple, $\sqrt{2}^2 = 2$. Il est moins
évident qu'on puisse élever un nombre irrationnel à une puissance
\emph{irrationnelle} et obtenir encore un résultat rationnel. Le
théorème suivant répond à cette question par l'affirmative.

\begin{thm}
Il existe des nombres irrationnels $a$ et $b$ tels que $a^b$ soit rationnel.
\end{thm}

\begin{proof}
Considérons $\sqrt{2}^{\sqrt{2}}$. Si ce nombre est rationnel, nous
avons terminé : il suffit de prendre $a = b = \sqrt{2}$. Sinon,
il est irrationnel. Or
\[
(\sqrt{2}^{\sqrt{2}})^{\sqrt{2}} = \sqrt{2}^{\sqrt{2} \cdot
  \sqrt{2}} = \sqrt{2}^2 = 2,
\]
qui est rationnel. Dans ce cas, nous prenons donc pour $a$ le nombre
$\sqrt{2}^{\sqrt{2}}$, et pour $b$ le nombre~$\sqrt 2$.
\end{proof}

Est-ce une démonstration valide ? La plupart des mathématiciens
répondent oui. Mais, là encore, quelque chose nous laisse sur notre
faim : nous avons établi l'existence d'un couple de nombres réels
ayant une certaine propriété sans dire \emph{quel} couple convient.
On peut démontrer le même résultat en donnant le couple $a$, $b$
\emph{dans} la démonstration : prenons $a = \sqrt{3}$ et
$b = \log_3 4$. Alors
\[
a^b = \sqrt{3}^{\log_3 4} = 3^{1/2 \cdot \log_3 4} = (3^{\log_3
  4})^{1/2} = 4^{1/2}= 2,
\]
puisque $3^{\log_3 x} = x$.
%
%
%
%

La logique intuitionniste vise à rendre compte d'un raisonnement où
des démarches comme celle de la première démonstration sont exclues.
Pour démontrer l'existence d'un $x$ satisfaisant~$!A(x)$, il faut donner
un~$x$ précis et démontrer qu'il satisfait $!A$, comme dans la seconde
démonstration. De même, démontrer sans hypothèses que $!A$ ou $!B$
est vraie exige une démonstration de l'une des deux propositions.
Cette exigence constructive ne signifie pas que toute disjonction
déduite d'hypothèses arbitraires permette d'en déduire séparément
un de ses membres.

Formellement, on obtient la logique intuitionniste en restreignant
d'une certaine façon un système de !!{derivation} pour la logique
classique. Du point de vue mathématique, il s'agit simplement de
systèmes déductifs formels ; mais, comme nous l'avons indiqué,
ils visent à rendre compte d'un certain raisonnement mathématique.
On peut y voir le raisonnement qu'autorise une conception philosophique
des mathématiques, telle que l'intuitionnisme de Brouwer ; on peut
aussi y voir un raisonnement plus « concret » et plus satisfaisant,
dans l'esprit du constructivisme de Bishop. On peut enfin discuter
de l'adéquation entre la description formelle et sa motivation informelle.
Quelles que soient nos positions philosophiques, nous pouvons étudier
la logique intuitionniste comme une logique présentée formellement ;
pour diverses raisons, nombre de logiciens mathématiciens trouvent
cette étude intéressante.

\end{document}
